\subsection{E ⊗\_A F 关于正合序列的性质}

命题5. 设 E、E'、E'' 为右 A-模，F 为左 A-模，且

\[
\mathrm{E} ^ {\prime} \stackrel {u} {\longrightarrow} \mathrm{E} \stackrel {v} {\longrightarrow} \mathrm{E} ^ {\prime \prime} \longrightarrow 0\tag{14}
\]

是线性映射的一个正合列。记 \(\bar{u} = u \otimes 1_{\mathbb{F}}, \bar{v} = v \otimes 1_{\mathbb{F}}\) ，则序列

\[
\mathrm{E} ^ {\prime} \otimes_ {\mathrm{A}} \mathrm{F} \xrightarrow {\bar {u}} \mathrm{E} \otimes_ {\mathrm{A}} \mathrm{F} \xrightarrow {\bar {v}} \mathrm{E} ^ {\prime \prime} \otimes_ {\mathrm{A}} \mathrm{F} \longrightarrow 0\tag{15}
\]

是 Z-同态组成的正合列。

根据第 2 号，公式（5）， \(\bar{v} \circ \bar{u} = (v \circ u) \otimes 1_{\mathbf{F}} = 0\) ；像 \(\mathrm{H} = \bar{u} (\mathbf{E}' \otimes \mathbf{F})\) 包含于核 \(\mathbf{L} = \operatorname{Ker}(\bar{v})\) ；通过取商，我们从 \(\bar{v}\) 得到一个 \(\mathbf{Z}\) -线性映射 \(f\) ，它把余核 \(\mathbf{M} = (\mathbf{E} \otimes \mathbf{F}) / \mathrm{H}\) （ \(\bar{u}\) 的）映入 \(\mathbf{E}'' \otimes \mathbf{F}\) ；必须证明 \(f\) 是双射的，因此只需定义一个 \(\mathbf{Z}\) -线性映射 \(g: \mathbf{E}'' \otimes \mathbf{F} \to \mathbf{M}\) ，使得 \(g \circ f\) 和 \(f \circ g\) 都是恒等映射。

设 \(x'' \in E'', y \in F\) ；根据假设，存在 \(x \in E\) ，使得 \(v(x) = x''\) 。我们证明，如果 \(x_1, x_2\) 是 \(E\) 的两个元素，满足 \(v(x_1) = v(x_2) = x''\) ，且 \(\phi: E \otimes F \to M\) 是典范映射，那么 \(\phi(x_1 \otimes y) = \phi(x_2 \otimes y)\) 。只需证明若 \(v(x) = 0\) 那么 \(\phi(x \otimes y) = 0\) ，这由以下事实得出： \(x = u(x')\) （其中 \(x' \in E'\) ），因此 \(x \otimes y = u(x') \otimes y = \bar{u}(x' \otimes y) \in H\) 。若把 \((x'', y)\) 映到 \(\phi(x \otimes y)\) （对所有 \(x \in E\) 满足 \(v(x) = x''\) 而言它取唯一的值），便定义了一个从 \(E'' \times F\) 到 \(M\) 的映射；该映射是 Z-双线性的，并且满足条件 (1)（第 1 号），因为 \(v(x\lambda) = x''\lambda\) 和 \((x\lambda) \otimes y = x \otimes (\lambda y)\) （对 \(x \in E\) ）；因此存在一个 Z-线性映射 \(g\) （从 \(E'' \otimes F\) 到 \(M\) ），使得 \(g(x'' \otimes y) = \phi(x \otimes y)\) （对 \(y \in F, x \in E\) 和 \(x'' = v(x)\) ）。该定义进一步证明了 \(f \circ g\) 在

\(E'' \otimes F\) 中形如 \(x'' \otimes y\) 的元素上与恒等映射一致，因此 \(f \circ g\) 是 \(E'' \otimes F\) 的恒等映射；另一方面，对于 \(x \in E\) 和 \(y \in F\) ，根据定义有 \(f(\phi(x \otimes y)) = v(x) \otimes y\) ，因此 \(g(f(\phi(x \otimes y))) = \phi(x \otimes y)\) ，并且由于形如 \(\phi(x \otimes y)\) 的元素生成 M， \(g \circ f\) 是 M 的恒等映射。

推论. 设 F、F'、F'' 为左 A-模，E 为右 A-模，且

\[
\mathrm{F} ^ {\prime} \stackrel {s} {\longrightarrow} \mathrm{F} \stackrel {t} {\longrightarrow} \mathrm{F} ^ {\prime \prime} \longrightarrow 0\tag{16}
\]

是线性映射的一个正合列。记 \(\bar{s} = 1_{\mathbb{E}} \otimes s\) , \(\bar{t} = 1_{\mathbb{E}} \otimes t\) ，则 \(\mathbf{Z}\) -同态的序列

\[
\mathrm{E} \otimes_ {\mathrm{A}} \mathrm{F} ^ {\prime} \stackrel {\hat {s}} {\longrightarrow} \mathrm{E} \otimes_ {\mathrm{A}} \mathrm{F} \stackrel {\hat {t}} {\longrightarrow} \mathrm{E} \otimes_ {\mathrm{A}} \mathrm{F} ^ {\prime \prime} \longrightarrow 0\tag{17}
\]

是正合的。

当 E（相应地，F）被视为左（相应地，右）A \(^{0}\) -模时，F ⊗A \(^{0}\) E 与 E ⊗A F 等同起来，并且对于 F' ⊗A \(^{0}\) E 和 F'\,' ⊗A \(^{0}\) E 也有类似的等同（第 1 号，命题 1 的推论 2）；该推论随后立即由命题5得出。

注. 注意一般来说，若 \(E'\) 是右 A-模 E 的一个子模，且 \(j:E'\to E\) 是典范单射，则典范映射

\[
j \otimes 1 _ {\mathbf {F}}: \mathbf {E} ^ {\prime} \otimes \mathbf {F} \rightarrow \mathbf {E} \otimes \mathbf {F}
\]

不一定是单射。换言之，对于正合列

\[
0 \longrightarrow \mathrm{E} ^ {\prime} \stackrel {u} {\longrightarrow} \mathrm{E} \stackrel {v} {\longrightarrow} \mathrm{E} ^ {\prime \prime} \longrightarrow 0\tag{18}
\]

一般不能得出结论说序列

\[
0 \longrightarrow \mathrm{E} ^ {\prime} \otimes \mathrm{F} \stackrel {u} {\longrightarrow} \mathrm{E} \otimes \mathrm{F} \stackrel {v} {\longrightarrow} \mathrm{E} ^ {\prime \prime} \otimes \mathrm{F} \longrightarrow 0\tag{19}
\]

是正合的。

例如取A = Z，E = Z， \(E' = 2Z\) ，F = Z/2Z。由于 \(E'\) 与 E 同构， \(E' \otimes F\) 与 E \(\otimes\) F 同构，而后者本身又与 F 同构（第 4 号，命题 4）。但对所有 \(x' = 2x \in E'\) （其中 \(x \in E\) ）及所有 \(y \in F\) , \(j(x') \otimes y = (2x) \otimes y = x \otimes (2y) = 0\) ，因为 2y = 0，且 \(E' \otimes F\) 在 E \(\otimes\) F 中的典范像化为 0。

换言之，必须注意区分（对于 E 的一个子模 \(E'\) 和某个元素 \(x \in E'\) ）：元素 \(x \otimes y\) “在 \(E' \otimes F\) 中计算''与元素 \(x \otimes y\) “在 \(E \otimes F\) 中计算''（换句话说，即元素 \(j(x) \otimes y\) ）之间的差别。

我们将在后面以平坦模的名义研究这样的模F，即对于每个正合序列(18)，序列(19)都是正合的（交换代数，I，§2）。

命题6. 给定两个正合序列(14)和(16)，同态 \(v \otimes t: \mathbf{E} \otimes_{\mathbf{A}} \mathbf{F} \to \mathbf{E}'' \otimes_{\mathbf{A}} \mathbf{F}''\) 是满射且其核等于

\[
\operatorname{Im} (u \otimes \mathrm{l} _ {\mathrm{F}}) + \operatorname{Im} (\mathrm{l} _ {\mathrm{E}} \otimes s)
\]

现在 \(v \otimes t = (v \otimes 1_{\mathbb{F}''}) \circ (1_{\mathbb{E}} \otimes t)\) （第 2 号，公式（5）），因此 \(v \otimes t\) 是满射的，因为根据命题5及其推论，它是两个满射同态的复合。另一方面， \(z \in \mathbb{E} \otimes \mathbb{F}\) 属于 \(v \otimes t\) 的核，当且仅当 \((1_{\mathbb{E}} \otimes t) (z)\) 属于 \(v \otimes 1_{\mathbb{F}''}\) 的核，即（根据 (15)）属于下列映射的像：

\[
u \otimes 1 _ {F ^ {\prime \prime}}: \mathrm{E} ^ {\prime} \otimes \mathrm{F} ^ {\prime \prime} \rightarrow \mathrm{E} \otimes \mathrm{F} ^ {\prime \prime}.
\]

但作为同态 \(t: \mathbf{F} \to \mathbf{F}''\) 是满射，因此

\[
\mathbf {l} _ {\mathbf {E} ^ {\prime}} \otimes t: \mathbf {E} ^ {\prime} \otimes \mathbf {F} \rightarrow \mathbf {E} ^ {\prime} \otimes \mathbf {F} ^ {\prime \prime}
\]

由命题5的推论可知是满射的，因此关于 \(z\) 的条件归结为存在一个 \(a \in \mathbf{E}' \otimes \mathbf{F}\) ，使得

\[
(1 _ {\mathrm{E}} \otimes t) (z) = (u \otimes t) (a).
\]

设 \(b = z - (u \otimes 1_{\mathbb{F}})(a)\) ；于是 \((1_{\mathbb{E}} \otimes t)(b) = 0\) ，并且根据 (17)， \(b\) 属于 \(1_{\mathbb{E}} \otimes s\) 的像，这证明了该命题。

换言之：

推论1. 设 \(\mathbf{E}'\) 是右 A-模 \(\mathbf{E}\) 的一个子模， \(\mathbf{F}'\) 是左 A-模 \(\mathbf{F}\) 的一个子模，而 \(\operatorname{Im}(\mathbf{E}' \otimes_{\mathbf{A}} \mathbf{F})\) 和 \(\operatorname{Im}(\mathbf{E} \otimes_{\mathbf{A}} \mathbf{F}')\) 为子 \(\mathbf{Z}\) -模（ \(\mathbf{E} \otimes_{\mathbf{A}} \mathbf{F}\) 中的），即典范映射 \(\mathbf{E}' \otimes_{\mathbf{A}} \mathbf{F} \to \mathbf{E} \otimes_{\mathbf{A}} \mathbf{F}\) ，

\[
\mathbf {E} \otimes_ {\mathbf {A}} \mathbf {F} ^ {\prime} \rightarrow \mathbf {E} \otimes_ {\mathbf {A}} \mathbf {F}.
\]

则存在一个典范的Z-模同构

\[
\pi : (\mathrm{E} / \mathrm{E} ^ {\prime}) \otimes_ {\mathrm{A}} (\mathrm{F} / \mathrm{F} ^ {\prime}) \rightarrow (\mathrm{E} \otimes_ {\mathrm{A}} \mathrm{F}) / (\operatorname{Im} (\mathrm{E} ^ {\prime} \otimes_ {\mathrm{A}} \mathrm{F}) + \operatorname{Im} (\mathrm{E} \otimes_ {\mathrm{A}} \mathrm{F} ^ {\prime}))\tag{20}
\]

使得，对于 \(\xi \in \mathbf{E} / \mathbf{E}'\) , \(\eta \in \mathbf{F}\) , \(\pi (\xi \otimes \eta)\) 是由所有元素 \(x \otimes y \in \mathbf{E} \otimes_{\mathbf{A}} \mathbf{F}\) （其中 \(x \in \xi\) 且 \(y \in \eta\) ）组成的类。

注意当 E 是一个 \(((B_{i}^{\prime}); A, (C_{j}^{\prime}))\) -多重模，F 是一个 \((A, (B_{h}^{\prime\prime}); (C_{k}^{\prime\prime}))\) -多重模，且 \(E^{\prime}\) 和 \(F^{\prime}\) 分别是 E 和 F 的子多重模时，同构 (20) 是对于两边的 \(((B_{i}^{\prime}), (B_{h}^{\prime\prime}); (C_{j}^{\prime}), (C_{k}^{\prime\prime}))\) -多重模结构的同构（第 3 号）。

推论2. 设 \(\mathfrak{a}\) 是 \(\mathbf{A}\) 的一个右理想， \(\mathbf{F}\) 是一个左 A-模，而 \(\mathfrak{a}\mathbf{F}\) 是子 \(\mathbf{Z}\) -模（ \(\mathbf{F}\) 中的），由形如 \(\lambda x\) 的元素生成，其中 \(\lambda \in \mathfrak{a}\) 且 \(x \in \mathbf{F}\) 。于是存在一个典范的 \(\mathbf{Z}\) -模同构

\[
\pi : (\mathrm{A} / \mathfrak {a}) \otimes_ {\mathrm{A}} \mathrm{F} \rightarrow \mathrm{F} / \mathfrak {a} \mathrm{F}\tag{21}
\]

使得，对所有 \(\bar{\lambda} \in \mathbf{A} / \mathfrak{a}\) 以及所有 \(x \in \mathbf{F}\) , \(\pi(\bar{\lambda} \otimes x)\) 是以 \(\mathfrak{aF}\) 为模的 \(\lambda x\) 的同余类，其中 \(\lambda \in \bar{\lambda}\) 。

特别地，当 A = Z 时，可以看出对于每个整数 n 和每个 Z-模 F， \((\mathbf{Z}/n\mathbf{Z}) \otimes_{\mathbf{Z}} F\) 被典范地等同于商 Z-模 F/nF。

推论 3. 设 A 是一个交换环，a 是 A 的一个理想，E 和 F 是两个 A-模，使得 a 包含在 F 的零化子中。则 (A/a)-模 E ⊗\_A F 与 (E/aE) ⊗\_\{A/a\} F 典范同构。

F 与 \(E \otimes_{A} F\) 被 a 零化，因此具有典范的 (A/a)-模结构（§1，第 12 号），并且如果我们记 \(E' = aE\) ，则 \(\operatorname{Im}(E' \otimes_{A} F) = 0\) ；于是存在 (20) 型的典范同构，它把 \(E \otimes_{A} F\) 映到 \((\mathrm{E}/a\mathrm{E}) \otimes_{A} F\) ，而后者本身又与 \((\mathrm{E}/a\mathrm{E}) \otimes_{A/a} F\) 相同（第 3 号，命题 2 的推论）。

推论4. 设 \(\mathfrak{a}, \mathfrak{b}\) 是交换环 \(\mathbf{C}\) 中的两个理想；则 \(\mathbf{C}\) -模 \((\mathbf{C} / \mathfrak{a}) \otimes_{\mathbb{C}} (\mathbf{C} / \mathfrak{b})\) 典范同构于 \(\mathbf{C}/(\mathfrak{a} + \mathfrak{b})\) 。

